\documentclass{claudeeditorial}

\renewcommand{\DocKicker}{Showcase}
\renewcommand{\DocTitle}{Claude Editorial LaTeX}
\renewcommand{\DocSubtitle}{Catálogo visual de componentes para documentos cálidos, densos y legibles}
\renewcommand{\DocAuthor}{Equipo de documentación}
\renewcommand{\DocRole}{Sistema editorial}
\renewcommand{\DocDate}{5 de mayo de 2026}
\renewcommand{\DocRef}{SHOWCASE-001}
\renewcommand{\DocFooter}{claudeeditorial showcase}

\renewcommand{\InterviewQuestionLabel}{Pregunta}
\renewcommand{\InterviewAnswerLabel}{Respuesta}

\hypersetup{
  pdftitle={Claude Editorial LaTeX Showcase},
  pdfauthor={Equipo de documentación}
}

\fancyhead[L]{\UIFont\footnotesize\color{claudeWarmText}\DocKicker}
\fancyhead[R]{\UIFont\footnotesize\color{claudeStone}\DocRef}

\titleformat{\section}
  {\DisplayFont\Large\bfseries\color{claudeInk}}{}
  {0pt}{}[\vspace{0.38em}{\color{claudeBorder}\titlerule}]
\titlespacing*{\section}{0pt}{1.4em}{0.6em}

\setstretch{1.22}

\begin{document}
\BodyCopy
\CornerRibbon[claudeTerracotta]{Demo}

{\small\UIFont\color{claudeTerracotta}\DocKicker}\par\vspace{0.2em}
{\fontsize{31}{35}\selectfont\DisplayFont\color{claudeInk}\DocTitle\par}
\vspace{0.2em}
{\normalsize\UIFont\color{claudeWarmText}\DocSubtitle\par}
\vspace{0.45em}
{\UIFont\small\color{claudeStone}\DocAuthor\hspace{1em}\textperiodcentered\hspace{1em}\DocRole\hspace{1em}\textperiodcentered\hspace{1em}\DocDate\par}

\vspace{0.45em}
\ClaudeBadge{Claude-inspired}
\hspace{0.35em}\ClaudeBadge[claudeForest]{Warm editorial}
\hspace{0.35em}\ClaudeBadge[claudeWarmText]{LaTeX class}
\hspace{0.35em}\ClaudeBadge[claudeGold]{PDF-ready}

\vspace{0.55em}
\OrnamentRule

\begin{claudehero}
\ClaudeLead{Esta clase traslada a LaTeX un lenguaje visual de interfaz: tokens de color tipo CSS, superficies cálidas, bordes suaves, jerarquía serif/sans y componentes reutilizables para componer documentos de trabajo con presencia editorial.}
\end{claudehero}

\section{Métricas y tarjetas}

\begin{tcbraster}[raster columns=2,raster equal height=rows,raster column skip=10pt,
  enhanced,colback=claudeIvory,colframe=claudeBorderSoft,boxrule=0pt]
\BigStat[\IconArrow[pulseOK]\,+18\,\%]{Lectura}{64\,pp}{Documento largo con ritmo de revista técnica.}
\BigStat[\IconArrow[claudeStone]\,estable]{Sistema}{20+}{Macros y entornos listos para recombinar.}
\end{tcbraster}

\vspace{0.6em}

\begin{tcbraster}[raster columns=4,raster equal height=rows,raster column skip=8pt,
  enhanced,colback=claudeIvory,colframe=claudeBorderSoft,boxrule=0pt]
\KPICard{Lead}{1 página}{resumen}
\KPICard{Color}{12 tokens}{paleta}
\KPICard{Cajas}{10+}{tcolorbox}
\KPICard{Salida}{PDF}{xelatex}
\end{tcbraster}

\section{Iconos, severidad y etiquetas}

\begin{tabularx}{\linewidth}{@{}>{\UIFont\small\bfseries\color{claudeInk}}p{0.18\linewidth}Y@{}}
\toprule
{\color{claudeStone}\small Icono} & {\color{claudeStone}\small Uso recomendado} \\
\midrule
\IconCheck\ check & operación completada, evidencia validada o condición satisfecha \\
\IconCross\ cross & riesgo confirmado, claim no sostenible o condición incumplida \\
\IconWarn\ warning & atención requerida, revisión pendiente o cambio sensible \\
\IconArrow\ arrow & transición, flujo, siguiente paso o movimiento \\
\IconStar\ star & prioridad alta, nota destacada o criterio clave \\
\IconInfo\ info & explicación auxiliar o contexto adicional \\
\IconLock\ lock & contenido restringido, privado o sujeto a revisión \\
\IconDot\ dot & marcador simple de lista o punto de control \\
\bottomrule
\end{tabularx}

\vspace{0.7em}
\SevCritical{bloqueante}\quad
\SevHigh{alto impacto}\quad
\SevMedium{revisar}\quad
\SevLow{menor}\quad
\SevInfo{contexto}

\vspace{0.7em}
\Tag{LaTeX}\ \Tag[pulseIdentidad]{XeLaTeX}\ \Tag[pulseConformal]{Editorial}\ \Tag[claudeTerracotta]{Claude-like}\ \Tag[claudeForest]{Open template}

\section{Citas y glosario}

\begin{PressQuote}[Principio de diseño]
Un documento no debería parecer una interfaz pegada en papel. La traducción correcta es conservar la jerarquía, la temperatura visual y la economía de acentos, pero dejar que LaTeX haga lo que sabe hacer: composición tipográfica.
\end{PressQuote}

\DefnEntry{Token visual}{Valor reusable de color, espaciado o jerarquía que mantiene consistencia entre piezas.}
\DefnEntry{Superficie}{Caja o banda que agrupa contenido sin convertir toda la página en una colección de tarjetas.}
\DefnEntry{Acento}{Uso limitado de terracota para navegación visual, decisión o énfasis semántico.}
\DefnEntry{Cadencia editorial}{Alternancia de texto, respiración, tablas, citas y bloques para que el lector no se fatigue.}

\section{Comparación}

\begin{Comparison}{}{}
\CompCol{Antes}{%
\IconCross\ Documento gris, sin jerarquía.\par
\IconCross\ Tablas densas sin ritmo.\par
\IconCross\ Acentos usados como decoración.\par
\IconCross\ Portada desconectada del resto.}
\CompCol{Después}{%
\IconCheck\ Sistema de color coherente.\par
\IconCheck\ Componentes con función clara.\par
\IconCheck\ Acento reservado para sentido.\par
\IconCheck\ Portada, desarrollo y cierre unidos.}
\end{Comparison}

\clearpage

\SectionOpener{01}{Bloques editoriales}{Cajas, notas y decisiones}

\begin{tcbraster}[raster columns=2,raster equal height=rows,raster column skip=10pt]
\begin{claudenote}
\textbf{\UIFont Nota editorial.}

Sirve para contexto, antecedentes, matices o avisos suaves.
\end{claudenote}

\begin{claudequote}
\textbf{\UIFont Idea central.}

\itshape La sensación debe ser cálida y precisa: suficiente aire, poco ruido y un acento reconocible.
\end{claudequote}

\begin{claudedecision}{Decisión}
Bloque para conclusiones, criterios o recomendaciones. Tiene mayor peso visual y conviene usarlo con moderación.
\end{claudedecision}

\begin{claudechecklist}{Checklist}
\begin{itemize}
\item mantener títulos breves;
\item reservar terracota para énfasis real;
\item evitar cajas innecesarias;
\item revisar saltos de página.
\end{itemize}
\end{claudechecklist}
\end{tcbraster}

\SectionOpener{02}{Panel oscuro}{Cambio de atmósfera}

\begin{claudedarkpanel}{Resumen de alto peso}
{\DisplayFont\fontsize{23}{27}\selectfont\color{claudeIvory}Un documento largo necesita cambios de ritmo.\par}

\vspace{0.5em}
{\color{claudeWarmSilver}Este panel funciona para manifiestos breves, aperturas de capítulo, resúmenes ejecutivos densos o secciones que necesitan más gravedad visual.}

\vspace{0.8em}
\ClaudeBadge[claudeWarmSilver]{dark panel}
\hspace{0.35em}\ClaudeBadge[claudeWarmSilver]{chapter opener}
\hspace{0.35em}\ClaudeBadge[claudeWarmSilver]{high contrast}
\end{claudedarkpanel}

\SectionOpener{03}{Teoremas y pruebas}{Documentos técnicos con tono editorial}

\begin{definicion}[Componente editorial]
Un componente editorial es una unidad composable que combina función semántica, jerarquía visual y espaciado reproducible.
\end{definicion}

\begin{teorema}[Consistencia de tokens]
Si todos los componentes consumen la misma paleta y las mismas reglas de jerarquía, el documento conserva identidad visual aunque cambie de estructura.
\end{teorema}

\begin{claudeproof}
\textbf{Demostración (esbozo).} Las cajas, badges, tablas e iconos dependen de los mismos nombres de color y del mismo par tipográfico. Al variar el contenido, la relación visual permanece acotada por esos tokens. \qed
\end{claudeproof}

\SectionOpener{04}{Entrevista}{Formato pregunta-respuesta}

\EntrevistaQ{¿Cuándo conviene usar esta clase en vez de una plantilla corporativa convencional?}
\EntrevistaA{Cuando el documento necesita parecer cuidado, legible y con identidad editorial, pero no debe convertirse en una landing page ni en una presentación.}

\SectionOpener{05}{Aperturas y cierre}{ChapterMark, DropCap y firma}

\ChapterMark{I}{Del documento como interfaz tranquila}

\DropCap{U}{na plantilla editorial no intenta decorar cada párrafo. Su trabajo es marcar rutas de lectura: dónde empieza una idea, qué merece atención y qué decisión debe recordarse al final.}

\begin{claudeepigraph}
El buen diseño documental deja que el contenido parezca inevitable.
\end{claudeepigraph}

\begin{claudedecision}{Criterio de cierre}
Usa menos piezas de las disponibles. La clase es completa para que puedas elegir, no para que todo aparezca en cada documento.
\end{claudedecision}

\ClaudeSignature{\DocAuthor}{\DocRole}

\end{document}
